% ============================================================================
% Chapter 2, Section 2: Motivation & Context
% ============================================================================

\subsection{Motivation \& Context}\label{sec:motivation}

\subsubsection{From Chatbots to Agents}\label{sec:chatbots_to_agents}

The modern era of large language models traces its roots to the Transformer architecture~\citep{vaswani2017attention}, whose self-attention mechanism enabled the training of models on vast corpora of text and unlocked fluent, general-purpose language generation. For the average consumer, the release of ChatGPT in late 2022~\citep{openai2022chatgpt} marked a turning point: hundreds of millions of users suddenly had access to a system that could draft essays, answer technical questions, and write code---all within seconds. These early interactions, however, exposed fundamental limitations. Models hallucinated plausible-sounding but factually incorrect information with high confidence, struggled with multi-hop reasoning, and had no mechanism for verifying their own outputs or learning from mistakes within a session.

These shortcomings spurred two parallel lines of architectural innovation. The first targeted \emph{reasoning depth}. Chain-of-thought prompting~\citep{wei2022cot} demonstrated that simply instructing a model to ``think step by step'' could dramatically improve performance on arithmetic and logical tasks. This insight was later internalized at the training level: OpenAI's o1 model~\citep{openai2024o1} applied large-scale reinforcement learning to teach models to produce extended chains of reasoning before answering, achieving competitive-programming and PhD-level science performance. DeepSeek-R1~\citep{deepseek2025r1} showed that such reasoning capabilities could emerge from reinforcement learning alone, without supervised fine-tuning on curated reasoning traces. The second line of innovation targeted \emph{grounding}. Retrieval-Augmented Generation (RAG)~\citep{lewis2020rag} addressed hallucination by conditioning model outputs on retrieved documents, while tool-use frameworks enabled models to call external APIs, execute code, and query databases during inference. The ReAct paradigm~\citep{yao2022react} unified reasoning and acting into a single interleaved loop, laying the groundwork for agents that could plan and execute multi-step tasks.

At the same time, a growing body of evidence suggests that the era of straightforward performance gains through model scaling is reaching its limits. Frontier models are already trained on several trillion tokens---approaching the estimated stock of high-quality human-generated text~\citep{villalobos2024run}---and the computational cost of continued model-level scaling is becoming prohibitive, with today's frontier models reaching sizes in excess of several trillion tokens. \tersoo{cite!}. This has shifted the frontier of progress from \emph{larger models} to \emph{better scaffolding}: the systems, protocols, and architectural decisions that surround a foundation model and determine how effectively it can operate in the real world.

The emergence of standardized protocols has accelerated this shift. The Model Context Protocol (MCP)~\citep{anthropic2024mcp}, introduced by Anthropic and now stewarded under the Linux Foundation, provides a universal interface for connecting language models to external tools, data sources, and services. By standardizing how agents interact with their digital environment, MCP and similar initiatives have transformed language models from impressive question-answering systems into autonomous actors capable of navigating codebases, operating software tools, and executing multi-step workflows that unfold over minutes, hours, or even days. It is this transition from single-turn chatbots to persistent, environment-interacting agents that motivates our exploration of the long-context evaluation setting that modern AI systems face.

\subsubsection{The Long-Context Setting}\label{sec:long_context}

The current generation of frontier AI systems operate increasingly in long-horizon, autonomous settings. Gone are the days where a human overseer interacts with an LLM at every step; rather, agent harnesses like OpenAI's Codex CLI or Anthropic's Claude Code are trusted with completing work at the level of human engineers. Enhanced reasoning capabilities combined with their ability to plan and interact with their environment in real time through tool-calling have extended the active lifetime of user requests to the order of minutes, hours, and even days~\citep{anthropic2025visible}. In fact, an early Google Gemini 2.5 technical report observed agents maintaining coherent behavior across hundreds of hours of continuous interaction~\citep{gemini2p5report}.

With newly focused attention on the long-context setting of agent interaction, highly publicized experiments have tested how far autonomous systems can scale. Cursor's FastRender project deployed hundreds of concurrent AI agents running autonomously for a week, producing over one million lines of Rust code and claiming to have  built a functional web browser from scratch~\citep{cursor2026scalingagents}. While the result was an interesting case study in long-context multi-agent orchestration, many have pointed out its poor code quality, over reliance on existing, complex libraries, performance issues, and failure to meet basic web standards like the Acid3 test \tersoo{cite!}, indicating a more incomplete state than initially marketed. In a similar spirit, Anthropic researcher Nicholas Carlini orchestrated 16 instances of Claude Opus 4.6 to collaboratively build a C compiler over two weeks, producing 100,000 lines of Rust capable of compiling the Linux kernel and passing 99\% of GCC's torture test suite~\citep{carlini2026ccompiler}. Yet the project faced criticism centering on its incomplete functionality, due to the compilers reliance on external tools for assembly and linkage as well as poor code efficiency \tersoo{CITE!}. These projects highlight a pattern in the performance of modern agent systems: achieving consistent, high-quality results in the long-horizon setting remains an open problem even for the most capable models.

These efforts operate at an industrial scale, but long-context agents are also beginning to empower individual users. Commercial coding agents such as Claude Code~\citep{anthropic2025claudecode} and open-source platforms like OpenHands~\citep{wang2024openhands} navigate multi-file codebases over extended sessions, performing planning, tool use, and context management at scales that would have been infeasible a year ago. In their totality, these developments highlight the urgency of  benchmarks that natively evaluate both models and agent harnesses in the long-context, embodied setting. Without standardized evaluation, it remains difficult to attribute successes and failures to the underlying model, the scaffolding architecture, or the interaction between the two. Closing this gap requires benchmarks where agents must sustain coherent reasoning across thousands of sequential decisions, manage growing context windows, and integrate multimodal perception with strategic planning.

Yet existing benchmarks largely fail to stress these capabilities simultaneously. Standard LLM evaluations focus on short-context question answering, coding, or mathematical reasoning. Game benchmarks tend to isolate individual axes: imperfect-information games emphasize short-episode strategic reasoning~\citep{brown2018superhuman, brown2019superhuman}, while open-world environments like Minecraft test exploration without adversarial pressure. Recent orthogonality analysis via BenchPress~\citep{papailiopoulos2026benchpress} has confirmed that Pok\'{e}mon gameplay captures capabilities not predicted by the 49 benchmarks in their evaluation matrix, suggesting that this domain tests a genuinely distinct and underserved axis of model competence.

\subsubsection{Game AI Benchmarks}\label{sec:game_benchmarks}

The challenges  of sustained coherence, context management, and integration of multi-modal perception with planning and action are not new to AI research.  Benchmarks have historically served to catalyze these advancements in AI capability. The Arcade Learning Environment (ALE)~\citep{bellemare2013arcade} spurred a decade of reinforcement learning (RL) progress, most notably through DeepMind's Deep Q-Networks (DQN)~\citep{mnih2015human}, by providing a standardized suite of Atari 2600 games. MineRL~\citep{guss2019minerl} extended this paradigm to open-ended environments with long-horizon objectives in Minecraft, exposing just how far RL methods were from real-world competence and stimulating research on the scaffolding layer surrounding the underlying machine learning algorithm---a theme later amplified by Video Pre-Training (VPT)~\citep{baker2022video}, which demonstrated that large-scale imitation learning from unlabeled gameplay videos could bootstrap capable agents in the same environment. More recently, NeurIPS competitions like Neural MMO~\citep{liu2023neurips, suarez2023neural} and Lux AI~\citep{tao2024lux} have pushed the frontier on multi-agent cooperation and resource management, focusing less on the ability of agents to solve single, isolated tasks and more on understanding the emergent behaviors of multi-agent systems. In each case, the availability of a standardized evaluation framework preceded and enabled methodological progress, foreshadowing today's agent paradigms.

A consistent pattern emerges across these results: RL achieves superhuman performance in fully observable, deterministic settings, but this margin erodes in stochastic, partially observable environments~\citep{samuel1959}. LLM-based agents consistently lag behind specialist RL and search systems in game settings~\citep{kaggle2025gamearena, paglieri2024balrog}, particularly when long-horizon planning and spatial reasoning are required. Hanabi~\citep{bard2020hanabi} and FightLadder~\citep{li2024fightladder} have advanced imperfect-information and competitive evaluation, while NetHack~\citep{kuttler2020nethack} and BALROG~\citep{paglieri2024balrog} benchmark long-horizon reasoning for both RL and LLM agents.

However, careful analysis reveals that few of these benchmarks combine adversarial reasoning with long-horizon planning at the level demanded of todays AI systems. As we discuss next, Pok\'{e}mon RPGs occupy a distinctive niche in this landscape---and the recent wave of projects, suffixed with the "-plays Pok\'{e}mon" moniker, have made the case for standardization especially urgent.

\subsubsection{Agents in Pok\'{e}mon Games}\label{sec:pokemon_agents}

Pok\'{e}mon has emerged as a prominent testbed for evaluating frontier AI systems, driven by a convergence of properties that make it uniquely challenging. Its combination of planning under partial observability, adversarial gameplay, and non-linear narrative structure has presented a challenging but surmountable gaming experience for hundreds of millions of players over the last three decades~\citep{pokemon2025franchise}. The franchise's RPG titles, in particular, feature expansive, partially observable worlds that require effective long-horizon planning, strategy adaptation, and knowledge integration to progress through a structured narrative. With approximately $10^{564}$ possible battle states~\citep{karten2025pokeagent}, team building across 1,000+ species with diverse movesets and abilities, and a competitive metagame that evolves continuously, the games present a complexity and dynamism that exceeds most existing benchmarks.

With the widespread adoption of LLMs following the release of ChatGPT in late 2022~\citep{openai2022chatgpt}, a natural question emerged: can models trained on a substantial fraction of all human-generated text---and therefore equipped with considerable world knowledge about Pok\'{e}mon mechanics, strategies, and walkthroughs---actually \emph{play} the games? Several high-profile projects have demonstrated both the appeal and the difficulty of answering this question. Claude Plays Pok\'{e}mon, a public Twitch stream, demonstrated extended thinking capabilities over approximately 35,000 actions to complete a small section of Pok\'{e}mon Red, but exposed struggles with basic spatial reasoning, vision interpretation, and obstacle navigation~\citep{engadget2025claudeplayspokemon, anthropic2025visible}. The Gemini Plays Pok\'{e}mon (GPP) project used Google's Gemini 2.5 Pro to play through Pok\'{e}mon Blue, ultimately completing the game in approximately 406 hours of cumulative gameplay, though the system frequently entered repetitive loops as the context window grew north of 100K tokens~\citep{zhang2025geminiplayspokemon, gemini2p5report}. Later GPP runs with Gemini 3 Pro demonstrated substantial improvements: completing Pok\'{e}mon Crystal without losing a single battle, using half the turns and 60\% fewer tokens to reach the same milestones as its predecessor~\citep{zhang2025gemini3playspokemon}. OpenAI's GPT-5 completed Pok\'{e}mon Red in 6,470 steps~\citep{engadget2025gptplayspokemon}, establishing a speed benchmark.

Taken in their totality, these projects served to highlight Pok\'{e}mon's offering as a distinctive form of out-of-distribution evaluation. While extensive Pok\'{e}mon knowledge exists in pretraining corpora in the form of move types, damage formulas, and walkthrough guidance, translating that latent knowledge into effective multi-turn decision-making under partial observability is fundamentally different from recognition and retrieval tasks. Moreover, while these efforts further validated Pok\'{e}mon as a testbed for embodied AI, in particular, they remain fundamentally fragmented. Different projects tested a number of models across different games (Red, Blue, Crystal, Emerald) across a non-standardized suite of evaluation criteria. Critically, they employed a variety of harness architectures with varying levels of built-in assumptions about what access the underlying models had to environment information. While impressive, one is still left to wonder if Gemini 3 Pro's completion of Crystal was comparable to Claude's progress in Red? Did GPT-5's step count account for the same mechanics as the GPP runs? By conflating harness engineering with model capability, it became impossible to attribute success to the agent architecture, the underlying model, or hardcoded assumptions that simplified perception and navigation. As we've discussed, the importance of standardized evaluation in games as a catalyst for the advancement of AI systems is well established. What remains central to this progress, however, is the existence of a standard. Pok\'{e}mon demands---and now receives---similar standardization through this benchmark. The gaps it exposes remain far from closed and motivate the benchmark, harness, and algorithmic contributions of this thesis.


